% Imporditud: tools/import_eksperimendid.py
% Allikas: Eesti füüsikaolümpiaadi eksperimendi ülesannete
%   kogu (Rael Kalda, 2023), ülesanne Ü153, lahendus L153.
\setAuthor{Erkki Tempel}
\setRound{lõppvoor}
\setYear{2019}
\setNumber{G E1}
\setDifficulty{4}
\setTopic{geomeetriline optika}
\setVahendid{Kumerpeegel, joonlaud (30 cm)}
\setPunktid{10}
\setAllikas{Eksperimendi kogu Ü153, lahendus lk 113}
\setLahendusseis{toores}

\prob{Kumerpeegel}
Määrake kumerpeegli fookuskaugus.

\hint

\solu
Kumerpeegel tekitab kujutise peegli taha. Asetame joonlaua peegli ette nii, et peeglis tekkiv kujutis on täpselt peegli laiune. See võimaldab meil mõõta kujutise suurust. Peegel peab olema meist võimalikult kaugel, et vähendada parallaksist tingitud viga. Mõõdame joonlaua pikkuse L, peegli diameetri d, joonlaua kauguse peeglist a ning peegli keskpunti kõrguse peegli serva tasapinnst x. Sarnastest kolmnurkadest $\triangle$AOF ja $\triangle$CDF saame kirja panna seose.

AO OH + HF = CD DF

Teades, et AO = 0,5L, CD = 0,5d, OH = a, HF = f ning DF = f $-$ x, saame leida peegli fookuse f

da + Lx f=

L$-$d

Suurema peegli (must) korral f $\approx$ 7,7 cm ning väiksema (kollane) korral f $\approx$ 16 cm.
\probend
